\documentclass[10pt]{article}

\usepackage[margin=1in]{geometry}
\usepackage{amsmath,amssymb,amsfonts}
\usepackage{booktabs}
\usepackage{graphicx}
\usepackage{microtype}
\usepackage{multirow}
\usepackage{url}
\usepackage{xcolor}
\usepackage[hidelinks]{hyperref}

\newcommand{\todo}[1]{\textcolor{red}{[TODO: #1]}}
\newcommand{\state}{\mathbf{x}}
\newcommand{\obs}{\mathbf{z}}
\newcommand{\cov}{\mathbf{P}}
\newcommand{\kalman}{\mathcal{K}}
\newcommand{\repr}{\mathcal{R}}
\newcommand{\loss}{\mathcal{L}}
\newcommand{\bbox}{\mathbf{b}}

\title{Event-Conditioned Residual Kalman Forecasting for UAV Bounding Boxes}
\author{
  \todo{Author names and affiliations}
}
\date{}

\begin{document}
\maketitle

\begin{abstract}
We study short-horizon forecasting of UAV bounding boxes in the FRED event-camera dataset. A carefully tuned constant-velocity Kalman filter is already a strong forecasting baseline, especially after extending the state from center position to full center-size boxes. We therefore evaluate learned models as residual dynamics on top of this baseline rather than as unconstrained box predictors. The residual model predicts acceleration corrections from recent box history, filtered state features, and, in image-conditioned variants, one event-derived image representation. Because standard splits can contain motion-prior shortcuts, we also measure how well fitted future acceleration is predicted from anchor position and velocity, and introduce a greedy sample decorrelation procedure that reduces this predictability and can optionally penalize global mean acceleration bias. Preliminary results suggest that event-conditioned residuals can improve box forecasting on the standard split, while decorrelated subsets provide a stricter diagnostic for whether gains come from visual evidence or from split-specific motion priors. \todo{Replace with final quantitative headline and exact comparison to prior FRED numbers.}
\end{abstract}

\section{Introduction}

Forecasting the future image-plane location of an observed UAV is useful for event-camera tracking, interception, sensor scheduling, and downstream recognition. The setting is challenging because the target can be small, fast, and partially represented by sparse events rather than dense RGB frames. Pure kinematic extrapolation is stable but cannot react to visual evidence about maneuvers. Conversely, an unconstrained neural predictor can exploit dataset shortcuts and produce physically implausible rollouts.

We investigate a hybrid approach: a constant-velocity Kalman model supplies a strong kinematic baseline, while a learned residual model predicts acceleration corrections conditioned on event representations, optional RGB, and the recent box history. The study is organized around three questions:

\begin{enumerate}
  \item How strong is an optimized Kalman filter when its process and measurement uncertainty parameters are tuned on the training split?
  \item Which visual representations contribute forecast information beyond kinematics: CSTR2/CSTR3, XT, YT, dataset event frames, RGB, and spatial cutouts around the anchor box?
  \item Do observed gains persist after reducing correlations between future acceleration and simple motion priors such as image position, velocity, distance-to-center, and speed?
\end{enumerate}

Our planned contributions are:

\begin{itemize}
  \item a reproducible residual Kalman forecasting benchmark on FRED UAV tracks using normalized box states and event/RGB representations;
  \item an ablation study over representation combinations and spatial cutout policies;
  \item a correlation analysis showing when future acceleration can be predicted from non-visual motion priors; and
  \item a sample-level decorrelation method based on greedy removal under a ridge-linear predictability score.
\end{itemize}

\todo{Tighten after final results: claim only what the completed experiments support.}

\section{Related Work}

\paragraph{Object Motion Forecasting.}
Classical trajectory forecasting often starts from constant-position, constant-velocity, or constant-acceleration assumptions, then augments them with learned dynamics or scene context. The prior FRED forecasting setup reports center-trajectory metrics, while our experiments emphasize full bounding-box forecasting and compare against an optimized center-size Kalman baseline. \todo{Add exact FRED and prior-work citations.}

\paragraph{Kalman Filtering and Learned Residuals.}
Kalman filters provide a probabilistic state estimate under linear-Gaussian dynamics. Hybrid approaches combine analytic transitions with neural residuals, allowing learned components to model unobserved forces or target maneuvers while retaining stable state evolution. Our model keeps the transition structure explicit and learns acceleration residuals in the normalized box state.

\paragraph{Event-Based Vision.}
Event cameras encode asynchronous brightness changes and are commonly converted into frame-like representations for convolutional networks. We evaluate multiple rendered representations, including XT, YT, CSTR variants, and dataset-native event frames, to separate visual information from kinematic priors.

\paragraph{Dataset Bias and Decorrelation.}
Forecasting datasets can contain correlations between target position, velocity, and future acceleration. Such correlations can create optimistic results for models that exploit split-specific motion priors rather than visual evidence. We therefore report both standard-split results and results on a decorrelated sample subset.

\todo{Replace placeholder citations with target workshop bibliography entries.}

\section{Problem Setup}

\subsection{State and Forecast Target}

Each sample is anchored at label time $t_0$. The observed history contains boxes up to and including $t_0$, and the target is the sequence of future boxes after $t_0$. A box is represented in normalized center-size coordinates
\begin{equation}
  \bbox_t = [c^x_t, c^y_t, w_t, h_t]^\top \in [0,1]^4,
\end{equation}
where coordinates are normalized by the frame size. The model state augments the box with first-order velocity:
\begin{equation}
  \state_t =
  [c^x_t, c^y_t, w_t, h_t, v^x_t, v^y_t, v^w_t, v^h_t]^\top.
\end{equation}

The default configuration uses a $400$ ms history window and an $800$ ms forecast window. Samples are generated from track-aligned FRED labels in \texttt{Event\_YOLO} and \texttt{cleaned\_tracks.txt}. Tracks and forecast-window candidates may be filtered by a minimum duration. Event inputs use trailing render windows to avoid future leakage.

\subsection{Inputs}

Let $\repr$ denote the configured set of rendered or dataset-native representations. The planned ablations include:
\begin{itemize}
  \item \textbf{CSTR2/CSTR3}: compact event representations rendered over the trailing event window;
  \item \textbf{XT/YT}: time-space projections, including aliases such as \texttt{xt\_my} and \texttt{yt\_mx};
  \item \textbf{event frames}: dataset-native frames from \texttt{Event/Frames} when available;
  \item \textbf{RGB}: \texttt{rgb} or \texttt{padded\_rgb} frames aligned to the anchor time; and
  \item \textbf{spatial cutouts}: masks or fixed crops around the final history box.
\end{itemize}

For XT inputs, spatial cutouts affect the $x$ axis only while preserving the temporal axis. For YT inputs, cutouts affect the $y$ axis only. Other image-like inputs use both spatial axes.

\subsection{Metrics}

Predictions are evaluated against future boxes using:
\begin{align}
  \mathrm{ADE}_{center} &= \frac{1}{T}\sum_{k=1}^{T}
  \left\|\hat{\mathbf{c}}_{t_k} - \mathbf{c}_{t_k}\right\|_2, \\
  \mathrm{FDE}_{center} &= 
  \left\|\hat{\mathbf{c}}_{t_T} - \mathbf{c}_{t_T}\right\|_2,
\end{align}
reported in pixels. We also report bounding-box ADE/FDE in pixels and mean IoU over the forecast horizon.

\todo{State exact train/validation/test folders once the final split policy is frozen.}

\section{Method}

\subsection{Optimized Kalman Baseline}

The baseline is a constant-velocity Kalman filter over $\state_t$. For a time step $\Delta t$, the transition matrix is
\begin{equation}
  \mathbf{F}(\Delta t) =
  \begin{bmatrix}
    \mathbf{I}_4 & \Delta t \mathbf{I}_4 \\
    \mathbf{0}_4 & \mathbf{I}_4
  \end{bmatrix},
  \qquad
  \state_{t+\Delta t} = \mathbf{F}(\Delta t)\state_t + \boldsymbol{\epsilon}_q.
\end{equation}
The measurement is the observed normalized box:
\begin{equation}
  \obs_t = \mathbf{H}\state_t + \boldsymbol{\epsilon}_r,
  \qquad
  \mathbf{H} = [\mathbf{I}_4 \; \mathbf{0}_4].
\end{equation}
The filter is initialized fresh for each sample, filters only the observed history window, and then rolls forward without future measurements. Initial, process, and measurement standard deviations are optimized on the training split. The primary optimized baseline uses differentiable backpropagation through the filter with log-space standard-deviation parameters and a weighted objective such as
\begin{equation}
  J = \mathrm{FDE}_{center} + \lambda_{ade}\mathrm{ADE}_{center}
      - \lambda_{iou}\mathrm{mIoU}.
\end{equation}

\subsection{Last-Four Linear Extrapolator}

For comparison, we also report a paper-style linear extrapolator. It estimates velocity for each box channel by a least-squares line fit over the last four observed boxes and then rolls forward with constant velocity:
\begin{equation}
  \hat{\bbox}_{t_k} = \bbox_{t_0} + \hat{\mathbf{v}}(t_k - t_0).
\end{equation}

\subsection{Residual Kalman Forecaster}

The learned model predicts acceleration residuals rather than boxes directly. It begins from either the last-four linear-fit state or the final filtered Kalman state. The main experiments use the optimized Kalman filter as the source of filtered state features. At forecast step $k$, the residual head predicts
\begin{equation}
  \mathbf{a}_k = f_\theta(\phi_u, \phi_h, \state_k, \Delta t_k),
\end{equation}
where $\phi_u$ contains the optional image feature and selected Kalman filter-derived features, and $\phi_h$ encodes recent box history. If no image, filter-state, or covariance feature is enabled, $\phi_u$ is omitted. The rollout is
\begin{align}
  \bbox_{k+1} &= \bbox_k + \mathbf{v}_k \Delta t_k
    + \frac{1}{2}\mathbf{a}_k \Delta t_k^2, \\
  \mathbf{v}_{k+1} &= \mathbf{v}_k + \mathbf{a}_k \Delta t_k.
\end{align}
The predicted box channels are clamped to $[0,1]$. A configuration option can restrict residual prediction to center acceleration only, leaving width and height under constant-velocity dynamics.

\subsection{Network Structure}

In the reported image-conditioned runs, the model receives at most one image representation. That image is encoded by a ResNet-18 style CNN with three input channels and a $128$-dimensional pooled output. This feature is concatenated with any selected Kalman filter-state or covariance features and passed through a learned projection MLP:
\begin{equation}
  \phi_u = g_\theta\left([\mathrm{enc}(\mathbf{I}), \phi_K]\right),
\end{equation}
where $\phi_K$ denotes the enabled filtered-state or covariance features. In non-image runs, the same projection is applied to $\phi_K$ alone; if no image, filter-state, or covariance features are enabled, this branch is absent. The default projection hidden dimension is $256$ with one linear layer followed by ReLU.

The history encoder receives the last $N_h=12$ history samples. Each sample contributes four normalized box values and one relative timestamp, giving an input dimension of $12 \times 5$. In the leakage-reduced setting, positions are expressed relative to the final history box. The residual head receives the concatenation of $\phi_u$, the history feature, current $8$-D rollout state, and scalar $\Delta t$, and outputs either $4$ acceleration channels or $2$ center-only acceleration channels. The main residual experiments predict all four box acceleration channels so that size dynamics can also deviate from constant velocity.

Optional Kalman state and covariance features can be appended to the projected feature branch. Supported state subsets are full state, center position and velocity (\texttt{center\_full}), center position, center velocity, and all velocity channels. Selected center positions can be mapped from $[0,1]$ to $[-1,1]$, placing the frame center at zero. Covariance features can be omitted, diagonal-only, or a flattened selected covariance submatrix; the current main experiments omit covariance features.

\subsection{Linear Residual Baselines}

To test how much of the residual task can be solved from non-image features alone, we also use a linearized version of the residual model. Setting the feature-projection, history, and residual hidden-layer counts to zero bypasses the intermediate MLPs and fits a single linear acceleration head over the selected features, rollout state, and $\Delta t$. We evaluate variants using only filtered center position, center velocity, center position plus velocity, or all velocity channels. These models are useful diagnostics: if they improve over Kalman, then part of the residual signal is already available from simple motion priors.

\subsection{Spatial Cutouts}

Spatial cutouts test whether the model uses local target evidence or global scene/motion shortcuts. For each anchor, inputs can be masked or cropped around the final history box. A box-scaled cutout uses the current box size multiplied by a scale factor, while a fixed-pixel cutout returns tensors of a configured size, e.g. $256 \times 256$. The fill value outside the retained region is fixed.

\subsection{Correlation Analysis}

To expose motion-prior bias, we fit a constant-acceleration curve to the center positions over each history-plus-future sample window:
\begin{equation}
  \mathbf{c}(t) \approx \mathbf{c}_0 + \mathbf{v}_0(t-t_0)
  + \frac{1}{2}\mathbf{a}(t-t_0)^2.
\end{equation}
This yields anchor position $\mathbf{c}_0$, velocity $\mathbf{v}_0$, and future acceleration $\mathbf{a}=[a_x,a_y]^\top$. We then measure how predictable $\mathbf{a}$ is from motion-prior features:
\begin{equation}
  \mathbf{x}_{prior} =
  [c^x_0-0.5, c^y_0-0.5, v^x_0, v^y_0,
   \|\mathbf{c}_0-[0.5,0.5]\|_2, \|\mathbf{v}_0\|_2].
\end{equation}

\subsection{Greedy Decorrelation}

The decorrelation method drops samples after the normal dataset cap is applied. We optionally first choose a reproducible random pre-subset, which allows multiple runs to start from exactly the same candidate pool. For a candidate subset $S$, sufficient statistics are accumulated for standardized prior features $X_S$ and fitted accelerations $Y_S$. The score combines mean absolute Pearson correlation, ridge-linear predictability, and optionally a penalty on global mean acceleration:
\begin{equation}
  \mathrm{score}(S) =
  \alpha \cdot \mathrm{mean}\left(|\mathrm{corr}(X_S,Y_S)|\right)
  + \beta \cdot \mathrm{mean}\left(\max(R^2,0)\right)
  + \gamma \cdot \left\|\mathrm{mean}(Y_S)\right\|_2.
\end{equation}
At each greedy step, the algorithm samples a candidate set of removable examples and removes the one that most reduces the score, stopping when the target keep fraction is reached. We also support a random-selection mode as a simple size-matched control that skips this score. The default feature mode uses centered position, velocity, distance-to-center, and speed. In the main decorrelated experiments we keep $50\%$ of capped samples and use a small mean-acceleration weight, $\gamma=10^{-4}$, to mildly discourage a constant vertical acceleration bias without making it dominate the unitless correlation and $R^2$ terms.

The standalone script also supports a coarser track-level variant that selects complete tracks using the same score and per-track sufficient statistics.

\section{Experiments}

\subsection{Protocol}

We use the optimized Kalman filter as the primary baseline and compare it against residual models. Unless otherwise noted, models use a $400$ ms history window, an $800$ ms forecast horizon, normalized boxes, AdamW training with learning rate $10^{-3}$, weight decay $10^{-4}$, Smooth L1 loss, and gradient clipping at norm $10$. Development runs use a folder-level train-evaluation split with $20\%$ of the training folders held out. The final held-out test split is evaluated only after model selection.

The main image experiments use fixed cutouts, relative box-history features, frame-centered filtered center positions, optimized Kalman state features, and ResNet-18 encoders. We do not use Kalman covariance features in the main results. Because training time is limited, we report a compact set of representation comparisons rather than a full factorial sweep.

\subsection{Kalman and Linear Diagnostics}

The first stage establishes how much performance is available without images. We compare the last-four linear extrapolator, the optimized center-size Kalman filter, and linear residual heads over selected filtered state features. These linear residual models use no image inputs and no box-history MLP.

\begin{table}[t]
  \centering
  \caption{Non-image forecasting baselines and linear residual diagnostics. Lower is better for ADE/FDE; higher is better for mIoU.}
  \label{tab:linear-diagnostics}
  \resizebox{\textwidth}{!}{%
  \begin{tabular}{lcccccc}
    \toprule
    Method & Learned features & Size residuals & ADE-C $\downarrow$ & FDE-C $\downarrow$ & ADE-BB $\downarrow$ & mIoU $\uparrow$ \\
    \midrule
    Last-four CV & none & no & \todo{} & \todo{} & \todo{} & \todo{} \\
    Optimized Kalman & none & no & \todo{} & \todo{} & \todo{} & \todo{} \\
    Linear residual & center position & yes/no & \todo{} & \todo{} & \todo{} & \todo{} \\
    Linear residual & center velocity & yes/no & \todo{} & \todo{} & \todo{} & \todo{} \\
    Linear residual & center position+velocity & yes/no & \todo{} & \todo{} & \todo{} & \todo{} \\
    Linear residual & velocities & yes/no & \todo{} & \todo{} & \todo{} & \todo{} \\
    \bottomrule
  \end{tabular}
  }
\end{table}

Preliminary runs indicate that filtered position and velocity can still predict useful acceleration residuals, even after decorrelation, although the resulting center improvements do not always translate into better mIoU. We therefore treat these models as diagnostics for motion-prior leakage rather than as the main contribution.

\subsection{Motion-Prior Correlation}

We compute acceleration-field summaries on the training, train-evaluation, and held-out test splits. For each split, we report mean absolute Pearson correlation and ridge-linear $R^2$ from motion-prior features to fitted center acceleration. We also visualize acceleration as a function of image position and velocity bins to expose edge-to-center trends, speed-dependent deceleration, and global vertical bias.

\begin{table}[t]
  \centering
  \caption{Motion-prior predictability of fitted future acceleration before and after sample decorrelation.}
  \label{tab:correlation}
  \begin{tabular}{lccccc}
    \toprule
    Split & Decorr. & Samples & Mean $|\rho|$ $\downarrow$ & Mean $R^2$ $\downarrow$ & $\|\bar{\mathbf{a}}\|_2$ $\downarrow$ \\
    \midrule
    Train & no & \todo{} & \todo{} & \todo{} & \todo{} \\
    Train & yes & \todo{} & \todo{} & \todo{} & \todo{} \\
    Train-eval & no & \todo{} & \todo{} & \todo{} & \todo{} \\
    Train-eval & yes & \todo{} & \todo{} & \todo{} & \todo{} \\
    Test & no & \todo{} & \todo{} & \todo{} & \todo{} \\
    Test & yes & \todo{} & \todo{} & \todo{} & \todo{} \\
    \bottomrule
  \end{tabular}
\end{table}

\subsection{Standard-Split Image Results}

The second stage evaluates event-conditioned residuals on the standard split. Early experiments suggest that event-derived inputs improve box metrics beyond the optimized Kalman baseline on this split. Dataset-native event frames have been the strongest single image input in the raw split among the tested formats, while CSTR2, CSTR3, XT, and YT provide complementary views of motion. We have not yet completed a broad search over all combinations, so the final table will focus on the strongest single inputs and a small number of combinations.

\begin{table}[t]
  \centering
  \caption{Standard-split forecasting results before decorrelation. Lower is better for ADE/FDE; higher is better for mIoU.}
  \label{tab:standard-results}
  \resizebox{\textwidth}{!}{%
  \begin{tabular}{lcccccc}
    \toprule
    Method & Inputs & Cutout & ADE-C $\downarrow$ & FDE-C $\downarrow$ & ADE-BB $\downarrow$ & mIoU $\uparrow$ \\
    \midrule
    Optimized Kalman & boxes & none & \todo{} & \todo{} & \todo{} & \todo{} \\
    Residual & filter state only & none & \todo{} & \todo{} & \todo{} & \todo{} \\
    Residual & event frames & fixed & \todo{} & \todo{} & \todo{} & \todo{} \\
    Residual & CSTR2 & fixed & \todo{} & \todo{} & \todo{} & \todo{} \\
    Residual & CSTR3 & fixed & \todo{} & \todo{} & \todo{} & \todo{} \\
    Residual & XT+YT & fixed & \todo{} & \todo{} & \todo{} & \todo{} \\
    Residual & event/CSTR + XT/YT & fixed & \todo{} & \todo{} & \todo{} & \todo{} \\
    Residual & event/CSTR + XT/YT + RGB & fixed & \todo{} & \todo{} & \todo{} & \todo{} \\
    \bottomrule
  \end{tabular}
  }
\end{table}

\subsection{Decorrelated-Subset Results}

The third stage repeats the strongest standard-split settings on the decorrelated sample pools. The decorrelation procedure keeps $50\%$ of capped samples, uses the motion-prior feature mode, and applies a small mean-acceleration penalty. Preliminary results suggest that decorrelation reduces but does not eliminate residual predictability from position and velocity. On decorrelated samples, CSTR2 has so far appeared stronger than dataset-native event frames, possibly because its normalized timestamp structure exposes local motion gradients that binary-like frames do not. This explanation remains speculative and should be treated as a hypothesis for follow-up experiments.

\begin{table}[t]
  \centering
  \caption{Forecasting results on the decorrelated sample split.}
  \label{tab:decorrelated-results}
  \resizebox{\textwidth}{!}{%
  \begin{tabular}{lcccccc}
    \toprule
    Method & Inputs & Keep frac. & ADE-C $\downarrow$ & FDE-C $\downarrow$ & ADE-BB $\downarrow$ & mIoU $\uparrow$ \\
    \midrule
    Optimized Kalman & boxes & 0.5 & \todo{} & \todo{} & \todo{} & \todo{} \\
    Residual & filter state only & 0.5 & \todo{} & \todo{} & \todo{} & \todo{} \\
    Residual & event frames & 0.5 & \todo{} & \todo{} & \todo{} & \todo{} \\
    Residual & CSTR2 & 0.5 & \todo{} & \todo{} & \todo{} & \todo{} \\
    Residual & CSTR2 + XT/YT & 0.5 & \todo{} & \todo{} & \todo{} & \todo{} \\
    Residual & best available combo & 0.5 & \todo{} & \todo{} & \todo{} & \todo{} \\
    \bottomrule
  \end{tabular}
  }
\end{table}

\subsection{Comparison to Prior FRED Results}

We compare against the original FRED box metrics and the recent center-forecasting results from \todo{cite prior arXiv work}. Our optimized Kalman extension to width and height appears to improve over the original FRED box results on mIoU and box ADE/FDE, while center ADE/FDE remains competitive but does not clearly surpass the strongest reported center-only methods. \todo{Insert exact numbers and ensure metric definitions match prior work.}

\subsection{Ablations Deferred by Scope}

The following ablations are useful but not required for the core workshop claim: cutout size, temporal bin count, temporal crop length for XT/YT, smaller CNN backbones, covariance features, RGB contribution, and broader representation combinations. We report them as limitations or future work unless completed before submission.

\section{Discussion}

The main empirical risk is overinterpreting standard-split residual gains. A residual model can improve forecasting by reading local event evidence, but it can also learn that certain positions and velocities imply a likely future acceleration under the dataset collection protocol. The linear residual diagnostics and decorrelation experiments are intended to separate these explanations. If a linear model over filtered state improves substantially, then some part of the residual signal is available without images. If image-conditioned gains survive decorrelation and exceed the filter-state-only residual, then the evidence for visual conditioning is stronger.

The decorrelated subset should not be interpreted as a true out-of-distribution benchmark. It removes a specified class of linear motion-prior shortcuts after sample capping, and the greedy candidate search is approximate. Nonlinear shortcuts, remaining split-specific behavior, and visual context cues can persist. The optional mean-acceleration penalty addresses a global acceleration bias, such as a downward tendency, but choosing its weight is application dependent: for a horizontal camera it may be a dataset artifact, while for some deployment settings gravity-related image-plane acceleration may be a real prior.

The preliminary single-input representation results also require caution. Event frames perform well on the raw split, while CSTR2 appears stronger in early decorrelated runs. One possible explanation is that CSTR2 preserves local motion direction through timestamp gradients, whereas event frames are closer to binary occupancy. However, this has not yet been isolated by controlled ablations over temporal bin count, crop size, or backbone capacity.

The strongest claim is therefore not that one event representation is universally best. Instead, the results should be interpreted as evidence that (i) full-box optimized Kalman filtering is a strong and necessary baseline on FRED, (ii) single-input event-conditioned residuals can improve box forecasting under the standard protocol, and (iii) correlation analysis and decorrelation are necessary diagnostics before attributing those gains to visual evidence or OOD-ready dynamics.

\section{Conclusion}

We presented a residual Kalman approach for short-horizon UAV bounding-box forecasting on FRED. The method extends center-only Kalman trajectory forecasting to a full center-size box state, then learns acceleration residuals from box history, filtered state features, and, in image-conditioned variants, a single event-derived image input. We also introduced linear residual diagnostics and a decorrelation procedure that reduces acceleration predictability from position and velocity, with an optional term for global mean acceleration bias. The resulting evaluation separates full-box claims against the original FRED benchmark from center-only comparisons against RPM-modulated Kalman forecasting. \todo{Replace with final quantitative summary after tables are populated.}


\bibliographystyle{plain}
\bibliography{references}

\end{document}
